\section{Introduction}\label{sec:intro}

\IEEEPARstart{R}{esistive} RAM (ReRAM) crossbars compute a matrix--vector product where the weights are
stored. Each weight is held as the conductance of one or more resistive cells; applying the input as
row voltages makes each column current the dot product of the input with that column's
conductances~\cite{shafiee2016isaac,chi2016prime}. Weights never move between memory and compute, which
removes most of the data movement that dominates the energy of neural-network inference on conventional
hardware. The cost is accuracy: a cell cannot be programmed to an exact conductance, its reading
fluctuates, and some cells are permanently stuck at their highest or lowest conductance.

A cell can store one bit (single-level cell, SLC) or several bits (multi-level cell, MLC). MLC is denser,
but its conductance levels are closer together, so the same conductance error corrupts more of the stored
value. HyFlexPIM~\cite{song2025hyflexpim} combines the two for transformers: it factorises each weight
matrix with singular value decomposition (SVD), keeps the components with the largest loss gradients in
SLC and stores the rest in MLC (\cref{sec:background}). Its evaluation applies relative noise to whole
weights at a single noise level.

We revisit this design with errors modelled cell by cell, as they arise in hardware, and ask which
components should be protected, how many, and which hardware effects still break the scheme. We study
BERT-Base~\cite{devlin2019bert} and DistilBERT~\cite{sanh2019distilbert} on the SST-2 sentiment
task~\cite{socher2013sst,wang2019glue}, and repeat the placement and mapping study on MNLI.

\textbf{Our contributions.} We keep HyFlexPIM's SVD factorisation and hybrid layout unchanged and add the
following.
\begin{itemize}
\item \textbf{Cell-level evaluation} (\cref{sec:model}). HyFlexPIM perturbs whole weights at one noise
level. We build a verified crossbar model that adds noise to every cell of a differential pair and its bit
slices, sweep programming noise from 0 to 10\%, and add read noise, level-dependent MLC noise, three
patterns of stuck cells, a fixed ADC with bit-serial inputs and IR drop. Every configuration is compared
on the same simulated chips.

\item \textbf{Which components to protect, and why} (\cref{sec:placement,sec:discussion}). HyFlexPIM ranks
components by their loss gradient. We show that ranking by singular value is as good as a
Fisher-information estimate of each component's expected damage and better than the gradient at high
noise, and explain this through the per-column quantisation of the crossbar.

\item \textbf{Stuck cells as the limiting effect.} HyFlexPIM does not model stuck cells. We find that every
other effect leaves the magnitude hybrid within about a point of all-SLC, while stuck cells break every
placement.

\item \textbf{Fault-aware mapping for SVD layers} (\cref{sec:mapping}). The components of an SVD layer can
be stored in any order. We use this freedom to place components around a tested chip's stuck cells, by a
Fisher-weighted assignment of components to slots, with no extra hardware or retraining.

\item \textbf{Tier-aware training} (\cref{sec:training}). In HyFlexPIM, hardware noise is applied only at
inference, not during fine-tuning. We train on simulated chips, with each component under the errors of the tier (SLC or MLC) it
will be stored in. Together with mapping, this lets the hybrid chip match an all-SLC chip with mapping up
to 1\% scattered stuck cells, using 34\% fewer cells.
\end{itemize}

\Cref{sec:background} reviews background and related work, \cref{sec:model,sec:placement,%
sec:mapping,sec:training} describe the model and methods, \cref{sec:setup,sec:results,sec:mnli} the experiments,
and \cref{sec:discussion,sec:conclusion} discuss and conclude.
